%\newcounter{AlgorithmCounter}[chapter] % defines algorithm counter for chapter-level
\newcounter{AlgorithmCounter}[section]
%\renewcommand{\theAlgorithmCounter}{\thechapter .\arabic{AlgorithmCounter}} %defines appearance of the algorithm counter
\renewcommand{\theAlgorithmCounter}{\thesection.\arabic{AlgorithmCounter}}
\DeclareCaptionLabelFormat{algocaption}{Algoritmo \theAlgorithmCounter} % defines a new caption label as Algorithm x.y

\lstnewenvironment{algorithm}[2][] %defines the algorithm listing environment
{   
    \refstepcounter{AlgorithmCounter} %increments algorithm number
    \captionsetup{labelformat=algocaption,labelsep=colon,font={normalsize}} %defines the caption setup for: it ises label format as the declared caption label above and makes label and caption text to be separated by a ':'
    \lstset{ %this is the stype
        backgroundcolor = \color{white},
        mathescape=true,
        frame=tb,
        numberstyle=\small, 
        basicstyle=\footnotesize\rmfamily,
        numbers=left,
        keywordstyle=\color{black}\bfseries,
        keywords={, for, input, output, return, for, and, or, to, datatype, function, goto, in, if, else, foreach, while, begin, end, print, import} %add the keywords you want, or load a language as Rubens explains in his comment above.
        numbers=left,
        xleftmargin=.04\textwidth,
        linewidth=.99\textwidth,
        #2}
}{}% this is to add specific settings to an usage of this environment (for instance, the caption and referable label)

% Environment for algorithmic pseudocode with styled background.
% #1 = caption text
\newenvironment{pseudocode}{%
    \begin{tcolorbox}[
        colback=codebackground,
        colframe=codebackground,
        boxrule=0pt,
        arc=3pt,
        top=4pt, bottom=4pt, left=6pt, right=6pt,
    ]%
}{%
    \end{tcolorbox}%
}

\newcounter{checksumfootnote}

% Repository del progetto: unico punto in cui sono definiti URL e versione
% citati nel testo. I link usano l'hash completo; il testo ne mostra la forma
% breve, che non esce dal margine.
\newcommand{\repourl}{https://github.com/gaiaLaRocca/parallel_computing_project}
\newcommand{\repocommit}{b386d1234ba5294c85c2cfb3b3fb49f70766da0f}
\newcommand{\repocommitshort}{b386d12}
% Link alla home del repository e alla versione citata.
\newcommand{\repolink}{\href{\repourl}{\texttt{github.com/gaiaLaRocca/parallel\_computing\_project}}}
\newcommand{\repoversion}{\href{\repourl/tree/\repocommit}{\texttt{\repocommitshort}}}
% \repofile{percorso}: link al file nella versione citata, es.
% \repofile{mandelbrot/src/mandelbrot_omp.cpp}.
\newcommand{\repofile}[1]{\href{\repourl/blob/\repocommit/#1}{\texttt{\detokenize{#1}}}}
